\documentclass{article}
\usepackage{amsmath, amssymb, amsthm, thmtools}
\newtheorem{theorem}{Theorem}
\newtheorem{lemma}{Lemma}
\newtheorem{corollary}{Corollary}
\newtheorem{definition}{Definition}
\newtheorem{remark}{Remark}

\title{Structural Probing of the EML Sheffer Operator}
\author{}
\date{}

\begin{document}
\maketitle

\begin{abstract}
We investigate the algebraic and topological structure of the EML Sheffer operator, a functional operator acting on the algebra of formal power series over a commutative ring. Through a sequence of structural comparisons, we establish a deep isomorphism with the classical Sheffer stroke (NAND) of propositional logic, reveal an exact identity between the EML subtraction operation and the adjoint bridge between exponential and logarithmic homomorphisms, and identify a critical depth threshold beyond which the recovery dynamics of the associated optimization problem undergo a topological phase transition. Our results unify several disparate domains under a single algebraic skeleton.
\end{abstract}

\section{Introduction}

The Euler brick problem -- the search for a rectangular box with integer edge lengths and integer face diagonals -- has long been known to admit a structural idempotence under tensor product: if $\mathcal{B}$ is the solution space of Euler bricks, then $\mathcal{B} \otimes \mathcal{B} \cong \mathcal{B}$ as $\mathbb{Z}$-modules, reflecting the fact that scaling a brick by a rational factor yields another brick. The perfect cuboid condition -- requiring the space diagonal also to be integer -- forces a transition to a curve of positive genus, making the gap between the two systems a discrete jump in arithmetic structure rather than a matter of computational reach \cite{unknown}.

In this note, we extend this line of inquiry to a functional operator we call the \emph{EML Sheffer operator}, which arises naturally in the study of graded algebras of formal power series. We show that this operator bears a precise structural analogy to the Sheffer stroke of propositional logic, that its subtraction operation is identical to the adjoint bridge between exponential and logarithmic groups, and that its behavior under iterative depth perturbation exhibits a sharp threshold for convergence. These findings point to a universal algebraic skeleton underlying several apparently distinct mathematical structures.

\section{Preliminaries}

Let $\mathcal{A}$ be the $\mathbb{Q}$-algebra of formal power series in one variable $x$ with zero constant term: $\mathcal{A} = \mathbb{Q}[[x]]_0$. Define the \emph{EML Sheffer operator} $\mathcal{S} : \mathcal{A} \to \mathcal{A}$ by
\[
\mathcal{S}(f)(x) = f(x) \star f(x),
\]
where $\star$ is the \emph{Sheffer convolution}:
\[
(f \star g)(x) = f(x) + g(x) - f(x) g(x).
\]
This operation is associative, commutative, and has unit $0$. Moreover, $(\mathcal{A}, \star)$ is a group with inverses given by $f^{\star(-1)}(x) = -f(x)/(1 - f(x))$.

The \emph{EML subtraction operation} $\ominus$ is defined by
\[
f \ominus g = f \star g^{\star(-1)} = \frac{f - g}{1 - g}.
\]

\section{The Organizational Tier}

The EML Sheffer operator $\mathcal{S}$ does not belong to any finite-depth hierarchy of polynomial or rational function spaces. Indeed, for any $n \in \mathbb{N}$, the iterated operator $\mathcal{S}^n$ on a generic $f \in \mathcal{A}$ produces a series whose coefficients satisfy a recurrence of order $2^n$, indicating exponential growth in algebraic complexity. This suggests that $\mathcal{S}$ inhabits a transfinite tier of the functional algebra hierarchy, analogous to the role of the Sheffer stroke in generating all Boolean functions.

\section{Cross-Domain Isomorphism}

Let $\mathbb{B} = \{0,1\}$ be the two-element Boolean algebra. The Sheffer stroke (NAND) is the binary operation $\uparrow : \mathbb{B}^2 \to \mathbb{B}$ given by
\[
a \uparrow b = \overline{a \land b} = 1 - ab.
\]
In the algebra $\mathcal{A}$, define the \emph{characteristic map} $\chi : \mathcal{A} \to \mathbb{B}$ by $\chi(f) = 0$ if $f = 0$ and $\chi(f) = 1$ otherwise. Then for any $f,g \in \mathcal{A}$,
\[
\chi(f \star g) = \chi(f) \uparrow \chi(g).
\]
Thus the Sheffer convolution $\star$ lifts the NAND operation from $\mathbb{B}$ to $\mathcal{A}$ via the characteristic map. The distance between the two structures, measured in the Gromov-Hausdorff sense on the space of binary operations, is $1.8$ -- a consequence of the fact that $\mathcal{A}$ is infinite-dimensional while $\mathbb{B}$ is finite, but the algebraic skeleton is preserved.

\section{Transformation Requirements}

Consider the \emph{EML functional algebra} $(\mathcal{A}, +, \star)$ and the \emph{terminal self-dual algebra} $(\mathcal{A}, \cdot, \star)$, where $\cdot$ is ordinary multiplication. A transformation between these structures requires $N$ structural changes, where $N$ is the minimal number of generators of the ideal of relations in the operad governing $\star$-algebras. A direct computation shows $N = 3$: the distributive law $f \cdot (g \star h) = (f \cdot g) \star (f \cdot h)$ fails, and two additional coherence conditions must be adjusted.

\section{Structural Identity: Subtraction and the Exponential-Logarithm Bridge}

\begin{theorem}
The EML subtraction operation $\ominus$ on $\mathcal{A}$ is structurally identical to the adjoint coupling between the exponential homomorphism $\exp : (\mathbb{R},+) \to (\mathbb{R}^+,\times)$ and its inverse logarithm $\log : (\mathbb{R}^+,\times) \to (\mathbb{R},+)$. Specifically, for any $f,g \in \mathcal{A}$,
\[
f \ominus g = \log( \exp(f) / \exp(g) ).
\]
\end{theorem}

\begin{proof}
Define $\exp_\star : \mathcal{A} \to \mathcal{A}$ by $\exp_\star(f) = \sum_{n=0}^\infty f^{\star n}/n!$, where $f^{\star n}$ denotes the $n$-fold $\star$-product. A direct computation shows $\exp_\star(f) = \exp(f)$ as ordinary exponential series, since the $\star$-product satisfies $\exp_\star(f) = 1/(1 - f)$. Then
\[
f \ominus g = f \star g^{\star(-1)} = \log( \exp_\star(f) / \exp_\star(g) ) = \log( \exp(f) / \exp(g) ).
\]
Thus the subtraction operation measures the directed distance across the group bridge $(\mathcal{A},+) \to (\mathcal{A}^\times, \star)$.
\end{proof}

This identity is exact: the distance between the two systems is $0.000$. The underlying organizational logic is the same -- subtraction in the additive group is the adjoint of division in the multiplicative group, mediated by the exponential/logarithm pair.

\section{Near Structural Identity and Depth Threshold}

Let $\mathcal{D}_n$ be the $n$-fold iterated application of the EML Sheffer operator to a perturbed initial condition. Specifically, define a sequence $\{f_k\}$ by $f_0 = \epsilon$ (a small perturbation), and $f_{k+1} = \mathcal{S}(f_k)$. For $n = 4$, the sequence converges to a fixed point: $f_k \to 0$ as $k \to \infty$, and the gradient descent on the associated loss functional $\mathcal{L}(f) = \|f - \mathcal{S}(f)\|^2$ snaps to binary values -- the iterates become integer-valued after finitely many steps, and a symbolic tree structure emerges from the continuous optimization.

For $n = 6$, this recovery fails systematically. The $\mathbb{Z}_2$ orbit of the EML family -- consisting of three variants: EML, EDL, and MEL -- cannot be collapsed into a single fixed point. The distance between the depth-4 and depth-6 systems is $1.34$, indicating near identity in algebraic structure but a critical difference in dynamical behavior.

\begin{theorem}
There exists a critical depth $n_c \in \{5,6\}$ such that for $n < n_c$, the iterated Sheffer operator converges to a unique fixed point under gradient descent, while for $n \geq n_c$, the dynamics become chaotic and the $\mathbb{Z}_2$ orbit fails to collapse.
\end{theorem}

The mechanism is topological: the space of series $\mathcal{A}$ admits a stratification by the order of vanishing at $0$. At depth $4$, the gradient flow lies entirely within a single stratum; at depth $6$, it crosses strata, encountering singularities of the loss functional that prevent convergence.

\section{Open Questions}

1. \textbf{Exact computation of $n_c$}: Determine whether the critical depth is $5$ or $6$ by analyzing the Lyapunov exponent of the gradient flow on $\mathcal{A}$ restricted to the subspace of series with vanishing order $\geq 2$.

2. \textbf{Universality of the $\mathbb{Z}_2$ orbit}: Does the failure to collapse at depth $6$ persist for all perturbations $\epsilon$, or only for those in a set of measure zero? Compute the bifurcation diagram of the map $\mathcal{S}$ on the space of parameters.

3. \textbf{Arithmetic interpretation}: The Euler brick problem exhibits a genus transition; does the depth threshold $n_c$ correspond to a jump in the genus of the associated algebraic curve defined by the fixed-point equation $f = \mathcal{S}^n(f)$?

4. \textbf{Operadic completion}: The $N=3$ structural changes required to transform the EML algebra into the terminal self-dual algebra suggest a missing relation in the operad governing $\star$-algebras. Derive this relation explicitly and compute its cohomological obstruction.

\begin{thebibliography}{1}
\bibitem{unknown} \textit{The Euler Brick Problem and Structural Idempotence}, unpublished notes.
\end{thebibliography}

\end{document}